\documentclass[12 pt]{exam}
\usepackage[margin=0.75in]{geometry}
\usepackage{amsmath,amssymb,color,amsthm}
\usepackage{enumerate,graphicx,multicol}
\usepackage{wrapfig}
\usepackage{pgf,tikz, subcaption, hyperref}
\usetikzlibrary{arrows}
\printanswers
\newtheorem{claim}{Claim}
\newtheorem{lemma}{Lemma}

\pagestyle{head}                       % This causes the exam to only have headers, no footers.
%\runningheadrule                     % This causes the headings to have an underline.

%\firstpageheader{Name:\enspace\hbox to 4.5in{\hrulefill} Section:\enspace\hbox to 1.2in{\hrulefill}}{}{}
%\extraheadheight{.25in}

\begin{document}
\hrule
\begin{center}
  \huge{Math 2710 Problem Set 13: Ch 13}    %Title
\end{center}
\hrule
\vskip .2in


\begin{questions}

\question  Explain the precise definition of a limit in your own words. What is the role of $\delta$ and of $\varepsilon$?

\begin{solution}\\
$\lim_{x \to a} f(x) = L \iff  \forall \varepsilon > 0, \exists \delta > 0, \forall x:
0 < |x - a| < \delta \implies |f(x) - L| < \varepsilon$.

$\delta$ is how far $x$ can be from $a$ so that $f(x)$ is at most $\varepsilon$ away from $L$.
\end{solution}

\question  Let $f(x)=3x+2$. We will show that $\displaystyle{\lim_{x \to 1} f(x) =5}$

\begin{parts}

\part Fill in the blanks:

\begin{solution}\\
For every $\varepsilon > 0$ there exists $\delta > 0$, such that $|f(x) - 5| < \varepsilon$
whenever $0 < |x - 1| < \delta$. 
\end{solution}


\part  Let $\varepsilon=0.5$.  Find $\delta$ so that $|f(x)-5|<\varepsilon$ whenever $|x-1|<\delta.$ Illustrate with a graph. 

\begin{solution}\\

\end{solution}


\part  Let $\varepsilon=0.01$.  Find $\delta$ so that $|f(x)-5|<\varepsilon$ whenever $|x-1|<\delta.$ Illustrate with a graph.

\begin{solution}\\

\end{solution}


\part  Find $\delta$(in terms of $\varepsilon$) so that $|f(x)-5|<\varepsilon$ whenever $|x-1|<\delta.$ Illustrate with a graph. 

\begin{solution}\\
Let $\varepsilon > 0$.

Let $\delta = \frac{\varepsilon}{3}$.

Assume that $0 < |x-1| < \frac{\varepsilon}{3}$.

\begin{align*}
    |x-1| <& \frac{\varepsilon}{3} \\
    |3x-3| <& \varepsilon \\
    |3x+2 - 5| <& \varepsilon
\end{align*}

Then for every $\varepsilon > 0$, there exists some $\delta > 0$ such that
$0 < |x - 1| < \delta \implies |f(x) - 5| < \varepsilon$.
\end{solution}


\end{parts} 

\question Prove, using the $\delta, \epsilon$ defintion of limits, that: $\lim_{x \to 3} 4x+6=18.$

\begin{solution}

\end{solution}


\question  Prove, using the $\delta, \epsilon$ defintion of limits, that: $$\lim_{x \to a} mx+b=ma+b, \textrm{ for } m\neq 0$$

\begin{solution}\\

\end{solution}


\question Prove that $f(x)=2x+1$ is continuous for all $x \in \mathbb{R}$.

\begin{solution}\\

\end{solution}


\question Prove that the following limit does not exist
$$\lim_{x \to 0} \frac{1}{x^2}$$

\begin{solution}\\

\end{solution}



\question Give an example where $\lim_{x \to c}f(x)$ does not exist and $\lim_{x \to c}g(x)$ does not exist but $$\lim_{x \to c}f(x)-g(x) =7.$$ 

\begin{solution}\\

\end{solution}




\question Prove that $$\lim_{x \to \infty} \frac{1}{x^2}=0.$$

\begin{solution}\\

\end{solution}








\newpage

\section*{Challenge Problems} 
\question Prove that if $\lim_{x \to c}f(x)=A$ and $\lim_{x \to c}g(x)=B$ then $\lim_{x \to c}f(x)-g(x)=A-B$.

\begin{solution}\\

\end{solution}


\question Prove the squeeze theorem: Suppose $g(x) \leq f(x) \leq h(x)$ for all $x \in \mathbb{R}$ satisfying $0 < |x-c| < \delta$. If $\lim_{x \to c}g(x)=l=\lim_{x \to c}h(x)$, then $\lim_{x \to c} f(x) =L$.

\begin{solution}\\

\end{solution}


\question Show that $$\lim_{x \to 3} x^2 =9.$$

\begin{solution}

\end{solution}


\end{questions}
\end{document}
